\subsubsection{Aplikace OpenAQ}

\[DRAFT\]

Druhou aplikací je americká aplikace \textit{OpenAQ} z New Yorku. Tato aplikace se, jak je uvedeno v dokumentaci API, zaměřuje především na měření kvality vnějšiho ovzduší. Podporovné veličiny jsou uvedeny v tabulce č. \ref{tab:open-aq-veliciny} \parencite{API}. Z té je zřejmé, že aplikace vůbec nemyslí na měření veličin akustického či světelného komfortu. 

Data v aplikaci OpenAQ jsou poskytována dvěma způsoby. 
Prvním z nich je přístup přes REST API, který umožňuje získat data ve formátu JSON. Tento přistup je vhodné pro programatické zpracování dat. Pro používání API jsou dostupné sady vývojových nástrojů (SDK) pro skriptovací jazyky Python a R. \parencite{API}  

Druhým způsobem je pravidelný archiv dat, který je dostupný na platformě AWS v S3 bucketu. Data jsou zde sdílena měsíčně ve formátu komprimovaných CSV souborů pomocí gzip. \parencite{zotero-item-552}

\begin{table}[!tbp]
  \centering
  \footnotesize
  \setlength{\tabcolsep}{2pt}
  \renewcommand{\arraystretch}{1.1}
  \caption{Sledované veličiny aplikací OpenAQ}
  \label{tab:open-aq-veliciny}
  \begin{tabular}{lllL{4cm}}
    \toprule
    \textbf{Kategorie} & \textbf{Veličina} & \textbf{Jednotka} & \textbf{Poznámka}\\
    \midrule
    \multirow{14}{*}{IAQ} 
                       & PM$_{1.0}$ & $\mu$g/m$^3$, cm$^{-3}$ & ---\\
                       & PM$_{2.5}$ & $\mu$g/m$^3$, ppm, cm$^{-3}$ & ---\\
     & PM$_{4.0}$ & $\mu$g/m$^3$ & ---\\
     & PM$_{10}$ & $\mu$g/m$^3$, cm$^{-3}$ & ---\\
     & CO & $\mu$g/m$^3$, ppm, ppb & ---\\
     & CO$_2$ & ppm & ---\\
     & NO & $\mu$g/m$^3$, ppm, ppb & ---\\
     & NO$_2$ & $\mu$g/m$^3$, ppm, ppb & ---\\
     & NO$_x$ & $\mu$g/m$^3$, ppm, ppb & ---\\
     & SO$_2$ & $\mu$g/m$^3$, ppm, ppb & ---\\
     & O$_3$ & $\mu$g/m$^3$, ppm, ppb & ---\\
     & BC & $\mu$g/m$^3$, ng/m$^3$ & Black Carbon (BC) - částice uhlíku menší než PM$_{2.5}$.\\ 
     & CH$_4$ & ppm & Metan. \\ 
     & UFP & cm$^{-3}$ & Ultra Fine Particles (UFP), někdy označována jako PM$_{0.1}$ jsou částice o velikostech v řádu desítek nanometrů.\\ 
     & P & hPa, mb & ---\\
    \midrule
    \multirow{1}{*}{Akustický komfort} & --- & --- & ---\\
    \midrule
    \multirow{2}{*}{Tepelný komfort} & $T$ & $^\circ$C, $^\circ$F & ---\\
     & RH & \% & ---\\
     & $v$ & m/s & ---\\
    \midrule
    \multirow{1}{*}{Světelný komfort} & --- & --- & --- \\
    \bottomrule
  \end{tabular}

  \footnotesize \textit{Zdroj:} Vlastní zpracování na základě dat z \parencite{API, zotero-item-549, UltrajemneCastice2025, CernyUhlik2025}. 
\end{table}

Za pomocí API byly identifikovány senzory, které jsou na plarformě OpenAQ použávány. Jejich shrnutí je uvedeno v tabulce č. \ref{tab04:openaq-supported-sensors}. V tabulce jsou uvedeny pouze veličiny, které jsou relevantní pro IEQ a byly v rámci rešerše autorem identifikovány. Je nutné poznamenat, že aplikace není závislá na konkrétním hardwaru. Uvedené senzory jsou pouze příklady, kdy autoři dat publikovali, jaký konkrétní hardware používají. To umožňuje uživatelům dat lépe porozumět, s jakou spolehlivostí a přesností data interpretovat. \parencite{API}

Praxe, kdy aplikace není závislá na konkrétním hardwaru je výhodná z pohledu udržitelnosti aplikace, kdy není potřeba implementovat podporu pro každý nový senzor. Tímto způsobem je břemeno napojení přeneseno na uživatele, který se rozhodne data publikovat. Tím se může vstupní bariéra pro publikování dat zvýšit, ale zároveň se zvyšuje flexibilita a škálovatelnost platformy.

{
  \footnotesize
  \setlength{\tabcolsep}{2pt}
  \renewcommand{\arraystretch}{1.1}
  \begin{longtable}{L{3.3cm}@{\hspace{8pt}}*{8}{c}@{\hspace{8pt}}*{3}{c}@{\hspace{8pt}}*{2}{c}@{\hspace{8pt}}*{2}{c}}
    \caption{Podporované instrumenty OpenAQ a jejich měřené veličiny relevantní pro IEQ}
    \label{tab04:openaq-supported-sensors}\\
    \toprule
    \multirow{3}{*}{\textbf{Senzor}} & \multicolumn{15}{c}{\textbf{Kategorie IEQ}} \\
                                     & \multicolumn{8}{c}{\textbf{IAQ}} & \multicolumn{3}{c}{\textbf{TK}} & \multicolumn{2}{c}{\textbf{SK}} & \multicolumn{2}{c}{\textbf{AK}} \\
                                     & PM\textsubscript{2.5} & PM\textsubscript{10} & O\textsubscript{3} & NO\textsubscript{2} & SO\textsubscript{2} & CO & CO\textsubscript{2} & TVOC & $T$ & RH & $v$ & $E$ & CCT & $L_A$ & $L$ \\
                                     \hline
Government Monitor & \cmark & \cmark & \cmark & \cmark & \cmark & \cmark & & & & & & & & & \\
\hline
Clarity Node-S (Clarity Sensor) & \cmark & & & \cmark & & & & & & & & & & & \\
\hline
HabitatMap AirBeam (HabitatMap Sensor) & \cmark & & & & & & & & & & & & & & \\
\hline
Senstate Sensor & & & & & & & & & & & & & & & \\
\hline
BEACO2N Sensor & \cmark & & & \cmark & & & \cmark & & \cmark & \cmark & & & & & \\
\hline
RAMP Sensor & & & \cmark & \cmark & \cmark & \cmark & & & & & & & & & \\
\hline
Aerosol Black Carbon Detector (ABCD) & & & & & & & & & & & & & & & \\
\hline
MA350 (AethLabs) & & & & & & & & & & & & & & & \\
\hline
Met One AIO 2 & \cmark & \cmark & & & & & & & & & & & & & \\
\hline
Met One BAM 1020 & \cmark & \cmark & & & & & & & & & & & & & \\
\hline
Met One BAM 1022 & \cmark & \cmark & & & & & & & & & & & & & \\
\hline
Serinus 30 (Ecotech) & & & & & & \cmark & & & & & & & & & \\
\hline
AirVisual (IQAir) & \cmark & & & & & & & & & & & & & & \\
\hline
AirGradient ONE (8PSL) & \cmark & \cmark & & & & & \cmark & \cmark & \cmark & \cmark & & & & & \\
\hline
AirGradient ONE (8PSL-DE) & \cmark & \cmark & & & & & \cmark & \cmark & \cmark & \cmark & & & & & \\
\hline
AirGradient ONE Gen 9 (I-9PSL) & \cmark & \cmark & & & & & \cmark & \cmark & \cmark & \cmark & & & & & \\
\hline
AirGradient ONE Gen 9 (I-9PSL-DE) & \cmark & \cmark & & & & & \cmark & \cmark & \cmark & \cmark & & & & & \\
\hline
AirGradient Open Air (O-1PST-CE) & \cmark & \cmark & & & & & & & \cmark & \cmark & & & & & \\
\hline
AirGradient Open Air Gen 1 (O-1PST) & \cmark & \cmark & & & & & & & \cmark & \cmark & & & & & \\
\hline
AirGradient Open Air (O-1PPT-CE) & \cmark & \cmark & & & & & & & \cmark & \cmark & & & & & \\
\hline
AirGradient Open Air Max Gen 1 (O-M-1PPST-CE) & \cmark & \cmark & & & & & & & \cmark & \cmark & & & & & \\
\hline
AirGradient Open Air Max Gen 1 (O-M-1PPSTON-CE) & \cmark & \cmark & & \cmark & & & & & \cmark & \cmark & & & & & \\
\hline
AirGradient DIY-Basic Gen 1-3 (I-1DIY) & \cmark & \cmark & & & & & \cmark & \cmark & \cmark & \cmark & & & & & \\
\hline
Aernode (Quanta S.r.l.) & \cmark & \cmark & \cmark & \cmark & \cmark & \cmark & & & \cmark & \cmark & & & & & \\
\hline
Miri Air (Miri Africa) & \cmark & \cmark & & & & & \cmark & & \cmark & \cmark & & & & & \\
    \bottomrule
  \end{longtable}

  \footnotesize \textit{Zdroj: }Vlastní zpracování za pomocí umělé inteligence na základě dat \parencite{API} a veřejně dostupných technických podkladů výrobců.

  \footnotesize \textit{Poznámka:} Seznam instrumentů OpenAQ obsahuje 27 položek. V tabulce jsou uvedeny jednotlivé instrumenty (bez slučování do rodin). Položky \textit{N/A} a \textit{Unknown AirGradient Sensor} nejsou uvedeny, protože nejde o konkrétní dohledatelná zařízení. Veličiny Black Carbon (např. \textit{MA350}, \textit{ABCD}) nejsou vyznačeny, protože nejsou samostatným sloupcem v použité IEQ matici.
}

Webové rozhraní aplikace je pro zobrazování dat na světové mapě podobné jako u aplikace Sensor.Community. Uživatelské rozrhraní je pro zobrazování dat přístupné bez předchozí regitrace. Uživatelé mohou zobrazení výsledků filtrovat podle: 
\begin{itemize}
  \item sledované znečisťující látky; a to:
    \begin{itemize}
      \item jakékoliv znečišťující látky (bez omezení),
      \item PM$_{2.5}$,
      \item PM$_{10}$,
      \item O$_{3}$,
      \item NO$_{2}$,
      \item SO$_{2}$,
      \item CO,
    \end{itemize}
  \item typu sledované lokality; těmi jsou:
    \begin{itemize}
      \item referenční stanice,
      \item lokace senzoru vzduchu,
      \item lokace senzoru, který nevykazuje aktivity,
    \end{itemize}
  \item dle partnerských projektů,
  \item a dle poskytovatele dat (171 poskytovatelů v době zpracování této části).
\end{itemize}

Tyto filtry jsou zobrazeny na obrázku č. 

